\begin{frame}{``왜 고전 ML이 아니라 딥러닝인가'' — 2축 + 2숫자}
  \begin{columns}[T]
    \begin{column}{0.5\textwidth}
      \kb{(축 1) 데이터의 구조} — 정형인가, 원시(비정형)인가
      \begin{itemize}
        \item 정형(수치/범주 열) $\Rightarrow$ 기본 전제:
          ``GBDT와 비교하라''
        \item 비정형(픽셀/음파/토큰) $\Rightarrow$ ``신경망이 특징
          추출기''라는 주장이 성립할 \textit{가능}이 생김
      \end{itemize}
      \vskip0.5em
      \kb{(축 2) 데이터의 규모} — 특징이 \textbf{공유}되는가,
      크기만큼 비용이 오르는가
      \begin{itemize}
        \item Ch10.1 파라미터 수 공식이 ``왜 CNN''을 수학적으로 설명
      \end{itemize}
    \end{column}
    \begin{column}{0.5\textwidth}
      \kb{(숫자 1) GBDT 베이스라인} — ``신경망이 정형 상한을
      넘는가''
      \vskip0.5em
      \kb{(숫자 2) 학습 곡선의 train/val 격차} — 신경망이
      \textbf{데이터가 커질수록 val 곡선이 계속 상승}한다면
      (GBDT는 일찍 포화), ``규모가 커지면 DL에 유리하다''는
      \textbf{전망}까지 제시 가능
    \end{column}
  \end{columns}
  \vskip0.6em
  \begin{block}{네 가지를 보고서와 발표에 담으면}
    ``왜 딥러닝인가''는 막연한 의문이 아니라 \textbf{반박 가능한
    주장}이 된다 — 16.2의 peer-review가 이 주장을 \textit{심사}하는
    기준이기도 하다.
  \end{block}
\end{frame}
